% 2.2.7
% Ficheiro incluído por main.tex

\subsection[Integração do Winbind com NSS]{Integração do Winbind com NSS}

\esquema{2-2-7-nss}{O NSS passa a perguntar ao Winbind pelos utilizadores e grupos do AD}

\passo{Passo 1 - Editar a configuração NSS}

\begin{lstlisting}[style=bashstyle]
sudo cp -a /etc/nsswitch.conf /etc/nsswitch.conf.bak
sudo nano /etc/nsswitch.conf
\end{lstlisting}

Deixar as linhas \texttt{passwd} e \texttt{group} assim (a linha \texttt{shadow} fica como está):

\begin{lstlisting}[style=cfgstyle]
passwd:         files systemd winbind
group:          files systemd winbind
shadow:         files systemd
\end{lstlisting}

Se aparecer \texttt{sss} nestas linhas, retirar: neste cliente não se usa o SSSD.

\passo{Passo 2 - Validar a integração}

\begin{lstlisting}[style=bashstyle]
getent passwd professor.teste
getent passwd aluno.teste
id professor.teste
id aluno.teste
getent group grupo-professores
getent group grupo-alunos
\end{lstlisting}

\emph{Verificação esperada:}

\begin{itemize}
\item \texttt{getent passwd} devolve a mesma linha do \texttt{wbinfo -i} (2.2.6, Passo 3);
\item \texttt{id professor.teste} inclui \texttt{grupo-professores} e \texttt{id aluno.teste} inclui \texttt{grupo-alunos}, além de \texttt{domain users};
\item \texttt{getent group} devolve cada grupo com o seu GID (ex.: \texttt{grupo-professores:x:11101:}).
\end{itemize}

\problemas{\hyperref[ec-b2]{Erros Comuns - 2.2 Winbind}}
